\documentclass[11pt]{article}
\usepackage{microstyle}

\begin{document}

\lessonheader{4}{The compensated law of demand, the Slutsky matrix, and the limits of WARP}{}

\section{The compensated law of demand, and WARP}

The compensated law of demand is an ``if and only if'' statement: for demand that is
homogeneous of degree 0 and satisfies Walras' law, its two parts are equivalent, so
WARP says the same thing as
\[ (P'-P)\cdot\bigl(x(P',w')-x(P,w)\bigr)\le 0
   \qquad\text{whenever}\qquad w'=P'\cdot x(P,w), \]
with strict inequality whenever $x(P',w')\ne x(P,w)$. That is, when the price change is accompanied by the wealth compensation that keeps
the original bundle just affordable, the compensated change in demand points in the
direction opposite to the price change. What follows proves the implication
WARP $\Rightarrow$ the inequality.

\section{Proof: WARP implies the compensated law of demand}

\topic{The trivial case}
If $x(P',w')=x(P,w)$ then $x(P',w')-x(P,w)=0$, so the left-hand side is zero and the
inequality holds.

\topic{The case $x(P',w')\ne x(P,w)$}
By construction of the compensated wealth, $P'\cdot x(P,w)=w'$, and by Walras' law at
the new prices the agent spends all of it, $w'=P'\cdot x(P',w')$. Combining the two,
\[ P'\cdot\bigl(x(P',w')-x(P,w)\bigr)=0. \tag{$*$} \]
Now use WARP. Since $P'\cdot x(P,w)=w'$, the old bundle $x(P,w)$ was affordable when
$x(P',w')$ was chosen, so
\[ x(P',w')\ \text{is revealed preferred to}\ x(P,w). \]
The opposite revealed preference cannot also hold, for that would be a
self-contradiction. If $x(P',w')$ had been affordable at $(P,w)$, then --- $x(P,w)$
being the bundle chosen there --- $x(P,w)$ would be revealed preferred to
$x(P',w')$, and the two bundles would each be revealed preferred to the other. WARP
rules this out, so $x(P',w')$ was \emph{not} affordable at $(P,w)$:
\[ P\cdot x(P',w')>w . \]
Subtracting $w=P\cdot x(P,w)$ --- Walras' law at the original prices --- from this,
\[ P\cdot\bigl(x(P',w')-x(P,w)\bigr)>0 . \]
Finally, use $(*)$:
\[
\begin{aligned}
  (P'-P)\cdot\bigl(x(P',w')-x(P,w)\bigr)
  &=P'\cdot\bigl(x(P',w')-x(P,w)\bigr)-P\cdot\bigl(x(P',w')-x(P,w)\bigr)\\
  &=0-P\cdot\bigl(x(P',w')-x(P,w)\bigr)<0 ,
\end{aligned}
\]
which is the compensated law of demand, strict in this case.

\section{The marginal price change: the substitution effect}

Now consider a small change in the price $p_j$ alone, with the wealth compensation
applied at the same time. The question is: by how much does the demand for good $i$
respond to $p_j$ when wealth is adjusted so that the \emph{old} bundle stays just
affordable? That response is the substitution effect $S_{ij}$.

\topic{Step 1: fix the starting point, and write compensated demand as a function of $p'$}
Assume demand is differentiable. Start at prices $p$ and wealth $w$, where the
consumer buys $x(p,w)$. Treat
$(p,w)$, and therefore the old bundle $x(p,w)$, as \emph{fixed numbers}: they
describe where we start and do not move. The only variable is the new price vector
$p'$. Slutsky compensation sets wealth at the new prices to
\[ w'(p')=p'\cdot x(p,w)=\sum_{k=1}^{l}p_k'\,x_k(p,w), \]
so the compensated demand for good $i$ is a function of $p'$ alone,
\[ h_i(p')\equiv x_i\bigl(p',\,w'(p')\bigr)=x_i\bigl(p',\,p'\cdot x(p,w)\bigr). \]

\topic{Step 2: see where $p_j'$ enters}
In $x_i\bigl(p',\,w'(p')\bigr)$ the new price $p_j'$ enters in two places:
\begin{flatnum}
  \item \emph{directly}, as the $j$-th price argument of the demand function;
  \item \emph{through wealth}, since $w'(p')=\sum_k p_k'\,x_k(p,w)$. Because the old
        bundle $x(p,w)$ is a constant here, only the $k=j$ term depends on $p_j'$,
        so
        \[ \frac{\partial w'}{\partial p_j'}=x_j(p,w): \]
        raising $p_j'$ by one unit raises the compensation by exactly the amount of
        good $j$ in the old bundle.
\end{flatnum}

\topic{Step 3: the chain rule}
Adding the two channels,
\[
  \frac{\partial h_i(p')}{\partial p_j'}
  =\underbrace{\frac{\partial x_i}{\partial p_j'}(p',w')}_{\text{direct price effect}}
   +\underbrace{\frac{\partial x_i}{\partial w}(p',w')}_{\text{wealth effect}}\;
    \underbrace{\frac{\partial w'}{\partial p_j'}}_{=\,x_j(p,w)} ,
  \qquad w'=p'\cdot x(p,w).
\]
Here $\partial x_i/\partial p_j'$ is the partial derivative of the demand function
$x_i(\cdot,\cdot)$ with respect to its $j$-th price argument, holding its wealth
argument fixed, and $\partial x_i/\partial w$ its derivative with respect to the
wealth argument; both are evaluated at the new point $(p',w')$. Note that
$x_j(p,w)$ carries the \emph{old} prices: it comes from $\partial w'/\partial p_j'$,
and the old bundle does not move.

\topic{Step 4: evaluate at $p'=p$}
$S_{ij}$ is the effect of an \emph{infinitesimal} change, so the derivative is taken
at the starting point, $p'=p$. There the compensated wealth is
$w'=p\cdot x(p,w)=w$ by Walras' law, so the evaluation point $(p',w')$ is just
$(p,w)$, and the price argument written $p_j'$ is then $p_j$. This gives the
\emph{substitution effect} of a change in the price of good $j$ on good $i$:
\[
  S_{ij}\equiv
  \frac{\partial h_i(p')}{\partial p_j'}\bigg|_{p'=p}
  =\frac{\partial x_i(p,w)}{\partial p_j}
   +\frac{\partial x_i(p,w)}{\partial w}\,x_j(p,w).
\]

\topic{Reading the formula}
The ordinary (uncompensated) price effect $\partial x_i/\partial p_j$ mixes two
things: relative prices change, \emph{and} the consumer becomes poorer (by $x_j$ per
unit rise in $p_j$, to first order). The compensation hands back exactly that loss,
and the term $\frac{\partial x_i}{\partial w}\,x_j$ is the demand response to it.
What is left, $S_{ij}$, is the response to the change in relative prices alone.
Rearranged, this is the \emph{Slutsky equation}:
\[
  \underbrace{\frac{\partial x_i(p,w)}{\partial p_j}}_{\text{total effect}}
  =\underbrace{S_{ij}}_{\text{substitution effect}}
   \;\underbrace{-\;\frac{\partial x_i(p,w)}{\partial w}\,x_j(p,w)}_{\text{wealth effect}} .
\]

\section{The Slutsky matrix}

The $l\times l$ matrix that collects all the $S_{ij}$ is the \emph{substitution
matrix}, or \emph{Slutsky matrix}. What can be said about its structure? The clue is
the compensated law of demand, which in differential form reads
\[ \sum_{i=1}^{l} dp_i\,dx_i\le 0 . \]
The change in wealth that accompanies the price change is
\[ dw=x(P,w)\cdot dP=\sum_{j=1}^{l}x_j(P,w)\,dP_j , \]
so the total differential of the demand for good $i$ is
\[
\begin{aligned}
  dx_i&=\sum_{j=1}^{l}\frac{\partial x_i(P,w)}{\partial P_j}\,dP_j
        +\frac{\partial x_i(P,w)}{\partial w}\,dw\\[2pt]
      &=\sum_{j=1}^{l}\frac{\partial x_i(P,w)}{\partial P_j}\,dP_j
        +\frac{\partial x_i(P,w)}{\partial w}\sum_{j=1}^{l}x_j(P,w)\,dP_j\\[2pt]
      &=\sum_{j=1}^{l}\Bigl[\frac{\partial x_i(P,w)}{\partial P_j}
        +\frac{\partial x_i(P,w)}{\partial w}\,x_j(P,w)\Bigr]dP_j\\[2pt]
      &=\sum_{j=1}^{l}S_{ij}\,dP_j .
\end{aligned}
\]

\section{Negative semi-definiteness of the Slutsky matrix}

Substituting $dx_i=\sum_j S_{ij}\,dP_j$ into the compensated law of demand,
\[
  \sum_{i=1}^{l}dP_i\,dx_i
  =\sum_{i=1}^{l}dP_i\sum_{j=1}^{l}S_{ij}\,dP_j
  =\sum_{i=1}^{l}\sum_{j=1}^{l}dP_i\,S_{ij}\,dP_j\le 0 ,
\]
by the compensated law of demand. Since this holds for every price change $dP$, the
quadratic form $dP' S\,dP$ is non-positive for every $dP$: the Slutsky matrix is
\emph{negative semi-definite}, in the sense $v'Sv\le0$ for all $v\in\R^{l}$, which
does not require $S$ to be symmetric.

Later in the course the same conclusion is obtained when axioms are imposed on
preferences. (The quadratic form only pins down the symmetric part
$\tfrac12(S+S')$. WARP does not deliver symmetry of $S$ in general --- with two goods
$S$ is automatically symmetric, given homogeneity and Walras' law, but not with three
or more --- whereas the preference-based derivation also gives symmetry.)

$S$ is never negative \emph{definite}. By Walras' law and then homogeneity of degree 0 (the Euler
equation of Lesson 3),
\[ \sum_{j}S_{ij}p_j=\sum_{j}\frac{\partial x_i}{\partial p_j}p_j
   +\frac{\partial x_i}{\partial w}\,p\cdot x(p,w)
   =\sum_{j}\frac{\partial x_i}{\partial p_j}p_j+\frac{\partial x_i}{\partial w}\,w=0, \]
so $S\,p=0$ and $p'Sp=0$: a proportional change in all prices, fully compensated,
changes nothing.

\section{WARP alone is not enough: an example}

Consider three price situations and the bundles chosen in each:
\[
\begin{array}{lll}
  P^{0}=(1,1,2), & P^{1}=(1,1,1), & P^{2}=(1,2,1),\\[2pt]
  x^{0}=(5,19,9), & x^{1}=(12,12,12), & x^{2}=(27,11,1).
\end{array}
\]
Here $x^{i}=x(P^{i},w^{i})$, so the wealth in situation $i$ is what $x^{i}$ costs at
$P^{i}$. The cost of each bundle at each price vector is

\begin{center}
\begin{tabular}{c|ccc}
        & $P^{0}$ & $P^{1}$ & $P^{2}$\\\hline
$x^{0}$ & \textbf{42} & 33 & 52\\
$x^{1}$ & 48 & \textbf{36} & 48\\
$x^{2}$ & 40 & 39 & \textbf{50}
\end{tabular}
\end{center}

\noindent
the bold entry in each column being the wealth, i.e.\ the cost of the bundle actually
chosen there: $w^{0}=42$, $w^{1}=36$, $w^{2}=50$. Reading the columns:

\begin{flatlist}
  \item At $P^{0}$ (wealth 42): $x^{1}$ costs 48 and is not affordable, but $x^{2}$
        costs 40. So $x^{2}$ was in the budget and $x^{0}$ was chosen:
        $x^{0}$ R.P.\ $x^{2}$.
  \item At $P^{1}$ (wealth 36): $x^{2}$ costs 39 and is not affordable, but $x^{0}$
        costs 33, so $x^{1}$ R.P.\ $x^{0}$.
  \item At $P^{2}$ (wealth 50): $x^{0}$ costs 52 and is not affordable, but $x^{1}$
        costs 48, so $x^{2}$ R.P.\ $x^{1}$.
\end{flatlist}

Chaining the three,
\[ x^{0}\ \text{R.P.}\ x^{2}\ \text{R.P.}\ x^{1}\ \text{R.P.}\ x^{0}, \]
a cycle: the revealed preferences are \emph{intransitive}. And yet WARP is satisfied
at every pair. In each of the three comparisons above the alternative bundle is
affordable only \emph{one} way round --- $x^{0}$ is not affordable at $P^{2}$, $x^{1}$
is not affordable at $P^{0}$, and $x^{2}$ is not affordable at $P^{1}$ --- so no two
bundles are ever revealed preferred to each other directly. WARP, which only ever
compares two situations at a time, never notices the cycle.

So WARP on its own is not enough to establish that the agent has rational (that is,
transitive) preferences. The example needs three goods: with only two goods, WARP
does rule out such cycles.

\section{The strong axiom of revealed preference (SARP)}

What is missing is transitivity. The \emph{strong axiom of revealed preference} adds
it: no chain of revealed preferences may return to its starting point,
\[ x^{0}\ \text{R.P.}\ x^{1}\ \text{R.P.}\ \cdots\ \text{R.P.}\ x^{0}
   \quad\text{is forbidden.} \]
SARP extends WARP from pairs to chains of any length (a chain of length two is
exactly what WARP forbids), so SARP implies WARP but not conversely; the example above
satisfies WARP and fails SARP, which is exactly why it is the counterexample it is.

\begin{transcriptnotes}
  \item The original opens with a title page reading only ``Lesson 4''.
  \item \textbf{Corrected:} in the total differential of $dx_i$ the original sums
        over $i$ inside the sum over $i$ and then keeps the index $i$ on the
        marginal effect of wealth,
        ``$\sum_j\frac{\partial x_i}{\partial P_j}dP_j
        +\frac{\partial x_i}{\partial w}\sum_ix_i(P,w)dP_i$''. The second sum must
        run over the good whose price changed, and the term multiplying
        $\frac{\partial x_i}{\partial w}$ is $x_j(P,w)\,dP_j$ --- which is what makes
        the bracket come out as $S_{ij}$ on the next line. Both places are written
        with $j$ here.
  \item \textbf{Corrected:} the derivative that defines the substitution effect is
        taken with respect to $p_j'$ (the price inside the compensated demand), not
        with respect to $p_j$; the notation $p'$ is used throughout.
  \item \textbf{Corrected:} the last page reads ``Strong axiom A.P.\ (SARP)''; the
        axiom is the strong axiom of \emph{revealed preference}. The line reads
        ``Assume away transitivity cycle'' (``away'' inserted above the line),
        i.e.\ SARP assumes away cycles of revealed preference --- the
        intransitive cycles of the example.
  \item Negative semi-definite: the notes conclude straight from
        $\sum_i\sum_j dP_iS_{ij}dP_j\le0$ that the Slutsky matrix is negative
        semi-definite, and that conclusion is correct as it stands, with
        semi-definiteness meant as $v'Sv\le0$ for all $v$ (as in MWG), which
        does not presuppose symmetry. \textbf{Corrected:} an earlier version of
        this transcription said the conclusion needed symmetry; the remark now says
        only that the quadratic form sees the symmetric part, and that symmetry is
        a separate property, obtained later from axioms on preferences.
  \item Page 5's remark ``Later, we get the same conclusion when we impose axioms
        on preferences'' is kept as the teacher's statement (it refers to negative
        semi-definiteness). \textbf{Corrected:} an earlier version of this
        transcription turned it into a claim that the \emph{next lecture} derives
        symmetry. \textbf{Added:} the parenthetical on symmetry.
  \item \textbf{Added:} the result $Sp=0$ (so $S$ is never negative definite), the
        assumptions (homogeneity, Walras' law) under which the ``if and only if''
        holds, the strict inequality when the bundles differ, and the remark that
        the WARP-without-SARP example needs three goods.
  \item \textbf{Corrected:} in the example, ``both were in the budget'' at $P^{0}$
        said $x^{1}$ was affordable there; it costs 48 against a wealth of 42, and
        only $x^{2}$ was in the budget.
  \item The chain-rule line wrote every partial derivative with the same composite
        argument $x_i(p',p'\cdot x(p,w))$, which makes the right-hand side look
        like the left-hand side. \textbf{Corrected:} an earlier version of this
        transcription wrote the direct term as $\partial x_i/\partial p_j$ while
        differentiating in $p_j'$; since the derivative is with respect to the
        new price, the direct term is $\partial x_i/\partial p_j'$, evaluated at
        $(p',w')$. It becomes $\partial x_i(p,w)/\partial p_j$ only after
        evaluating at $p'=p$. \textbf{Added:} the derivation is split into four
        steps (fix the starting point; the two channels through which $p_j'$
        enters; the chain rule; evaluation at $p'=p$ using Walras' law), with the
        compensated demand named $h_i(p')$, and the Slutsky equation is stated.
  \item ``slatsky'' is spelled \emph{Slutsky} here, and ``intransive'' is
        \emph{intransitive}.
  \item The original uses lower case $p$ in the substitution-effect section and upper
        case $P$ in the Slutsky-matrix section; the case is kept as written, as
        elsewhere in these notes.
  \item The example is transcribed with its numbers unchanged; the cost table was
        recomputed from the price vectors and bundles and matches the original entry
        by entry ($x^{0}$: 42, 33, 52; $x^{1}$: 48, 36, 48; $x^{2}$: 40, 39, 50).
        The original writes the table with the bundles as rows and the price vectors
        as columns, which is kept.
  \item The line that introduces the marginal-price-change section reads ``we move
        to key results readings marginal price change'' in the original; the word
        after ``results'' is probably ``regarding''. The sentence itself is not
        typeset; the section heading replaces it.
  \item \textbf{Corrected:} the original introduces the marginal change as ``change
        in $x_i$ when $P_i$ changes''; the derivation that follows (and $S_{ij}$)
        concerns a change in the price $p_j$ of good $j$, so $p_j$ is used.
  \item On the proof page, the original's step ``$x(P',w')$ is R.P.\ to $x(P,w)$, so
        WARP tells us that $x(P,w)$ was picked under $(P,w)$ because $x(P',w')$ was
        not affordable'' compresses the argument: that $x(P',w')$ was not affordable
        at $(P,w)$ is the \emph{conclusion} drawn from WARP (the opposite revealed
        preference would be a contradiction), not a premise. The proof is written out
        with that step spelled out.
  \item Pages of the original, for tracing back: page 2 the ``iff'' statement and the
        proof of WARP $\Rightarrow$ the inequality; page 3 the marginal price change
        and $S_{ij}$; page 4 the Slutsky matrix and the differential form; page 5 the
        negative semi-definiteness conclusion; page 6 the three-price example;
        page 7 the strong axiom.
\end{transcriptnotes}

\end{document}
